% Classic textbook style for the MIT 6.5130 lecture notes.
\usepackage[paperwidth=7in,paperheight=10in,inner=0.9in,outer=0.9in,top=0.95in,bottom=1in]{geometry}
\usepackage{newpxtext,newpxmath}
\usepackage{fontspec}
\setmonofont{Inconsolatazi4}[Extension=.otf, UprightFont=*-Regular, BoldFont=*-Bold, Scale=MatchLowercase, HyphenChar=None]
\usepackage{microtype}
\usepackage{graphicx}
\usepackage{amsmath}
\usepackage{xcolor}
\usepackage{listings}
\usepackage[numbers,sort&compress]{natbib}
\usepackage{titlesec}
\usepackage{fancyhdr}
\usepackage{emptypage}
\usepackage{adjustbox}
\usepackage{xurl}

% Code
\definecolor{codebg}{gray}{0.955}
\lstset{
  basicstyle=\ttfamily\small,
  columns=fullflexible, keepspaces=true,
  backgroundcolor=\color{codebg},
  frame=none, xleftmargin=1.2em, xrightmargin=1.2em,
  framexleftmargin=0.6em, framexrightmargin=0.6em,
  aboveskip=1.1ex, belowskip=1.1ex,
  breaklines=true, tabsize=4, upquote=true,
  extendedchars=true, escapeinside={(*@}{@*)},
  literate={->}{{$\rightarrow$}}2 {<=}{{$\leq$}}1 {>=}{{$\geq$}}1
           {λ}{{$\lambda$}}1 {→}{{$\rightarrow$}}1 {≤}{{$\leq$}}1 {≥}{{$\geq$}}1,
}
\newcommand{\code}[1]{\texttt{#1}}
% MathJax built-ins used by the notes.
\providecommand{\lt}{<}
\providecommand{\gt}{>}
\input{macros}
\newenvironment{aside}{\begin{quote}\small}{\end{quote}}

% Inline slide images: non-floating so sequences of frames can break across pages.
\newlength{\slidewidth}
\setlength{\slidewidth}{\linewidth}
\newenvironment{slidefig}{\par\medskip\begin{center}}{\end{center}\medskip}

% Display math that shrinks to the text width only when it would overflow.
\newcommand{\fitdisplay}[1]{\par\smallskip\begin{center}\adjustbox{max width=\linewidth}{$\displaystyle #1$}\end{center}\smallskip}

% Note for lectures whose page on the course site is not yet written.
\newcommand{\sitenote}[1]{\begin{quote}\itshape #1\end{quote}}

% Headings
\titleformat{\chapter}[display]
  {\normalfont\Large\scshape}{Lecture \thechapter}{1.2ex}{\normalfont\upshape\huge\bfseries\raggedright}
\titlespacing*{\chapter}{0pt}{0pt}{3.5ex}
\titleformat{\section}{\normalfont\large\bfseries}{\thesection}{0.8em}{}
\titleformat{\subsection}{\normalfont\normalsize\bfseries}{\thesubsection}{0.8em}{}
\titleformat{\part}[display]
  {\normalfont\centering\Large\scshape}{Unit \thepart}{2ex}{\normalfont\upshape\Huge\bfseries}

% Running heads
\pagestyle{fancy}\fancyhf{}
\fancyhead[LE]{\small\thepage\quad\scshape Introduction to Program Synthesis}
\fancyhead[RO]{\small\scshape\nouppercase{\leftmark}\quad\upshape\thepage}
\renewcommand{\chaptermark}[1]{\markboth{#1}{}}
\renewcommand{\sectionmark}[1]{}
\renewcommand{\headrulewidth}{0pt}
\fancypagestyle{plain}{\fancyhf{}\fancyfoot[C]{\small\thepage}}

\usepackage[hidelinks,bookmarksnumbered]{hyperref}
\renewcommand{\bibname}{References}
\setcounter{tocdepth}{1}
\setlength{\emergencystretch}{3em}
\raggedbottom
